\section*{Chapter \ref{ch:iv2:options-parity-and-greeks} --- Options: Parity and Greeks}
\markboth{\protect\omTocMark{}{Answers: Options: Parity and Greeks}}{\protect\omTocMark{}{Answers: Options: Parity and Greeks}}

\begin{solution}{iq:iv2:options-parity-and-greeks:1}
Your largest gain is \$7, your largest loss is \$3, and you break even at 103. It is a moderately bullish view: a bull call spread\index{call spread}.

The net cost is $5 - 2 = 3$. Below 100, both expire worthless, and you lose 3. Above 110, the spread pays the full 10, so you make $10 - 3 = 7$. In between, it pays $S - 100$, so you break even at 103.

\iqwrong{Thinking the short call adds risk. It caps your gain, and that is what pays for part of the call you bought.}
\end{solution}

\begin{solution}{iq:iv2:options-parity-and-greeks:2}
You break even at 92 and 108. At 105 you lose \$3, and at 120 you make \$12. You are betting that the stock moves a lot, in either direction, and more than the options' price implies.

The payoff is $|S - 100|$, and you paid 8. So at 105, $5 - 8 = -3$, and at 120, $20 - 8 = 12$. The largest loss is 8, at 100.

\iqwrong{Thinking any move wins. The stock must move more than the \$8 you paid, and a \$5 move still loses.}
\end{solution}

\begin{solution}{iq:iv2:options-parity-and-greeks:3}
Your largest profit is \$8, at 100. Your largest loss is the \$2 you paid. You break even at 92 and 108. It cannot have a negative price, because its payoff is never negative.

At 100, the 90 call pays 10, and the others pay nothing. So the payoff rises from 0 at 90 to 10 at 100, and it falls back to 0 at 110. Less the cost of 2, the profit is 8 at the top, and it is zero at 92 and 108.

A position that can only pay something must have a positive price. If not, someone pays you to take it, and this is called a butterfly\index{butterfly} arbitrage.

\iqwrong{Assuming the two short calls bring an unlimited loss. Above 110, the two long calls cancel them exactly.}
\end{solution}

\begin{solution}{iq:iv2:options-parity-and-greeks:4}
It is a long forward at 100, and it behaves like owning the stock. The stock must be at 102, because the position's cost is $S - K$ with zero rates. If the stock finishes at 110, you make \$8.

At expiry, the call pays $\max(S - 100, 0)$ and the short put pays $-\max(100 - S, 0)$. Together they pay $S - 100$. By put-call parity\index{put-call parity} with $r = 0$, $C - P = S - K$, so $S = 102$. The profit at 110 is $10 - 2 = 8$.

\iqwrong{Thinking this has limited risk because it is made of options. A short put has all the downside of the stock, and it is hiding here.}
\end{solution}

\begin{solution}{iq:iv2:options-parity-and-greeks:5}
The lowest price is $100 - 90e^{-0.05} \approx \$14.39$. At \$12, buy the call, short the stock and lend the proceeds. You make at least about \$2.39 in today's terms.

Today you pay 12 and receive 100, so you lend 88 at 5\%. That grows to $88e^{0.05} \approx 92.51$. At expiry, if $S > 90$, you exercise and buy at 90 to close the short, for a gain of about 2.51. If $S < 90$, you buy at $S$, for a gain of $92.51 - S$, which is more. Discounted, 2.51 is about 2.39.

\iqwrong{Using the bound $S - K = 10$, for which 12 looks fine. For a European call with positive rates, the strike is paid later, so it counts at its present value.}
\end{solution}

\begin{solution}{iq:iv2:options-parity-and-greeks:6}
It is too dear because the most a put can pay is 100, and that is one year away, so it is worth at most $100e^{-0.05} \approx \$95.12$ today. Sell it at \$97 and lend the money.

Lend 97 at 5\% for a year, which grows to about 101.97. The most you can owe is 100, if the stock goes to zero. So you make at least 1.97 at expiry, which is the 1.88 gap today.

\iqwrong{Comparing the put with the strike, 100, and calling \$97 fine. The strike is received at expiry.}
\end{solution}

\begin{solution}{iq:iv2:options-parity-and-greeks:7}
A put with a higher strike must cost at least as much, and here it costs less. Buy the 110 put and sell the 100 put, which pays you \$1 up front. Your worst outcome is the \$1 profit. The gap between the puts can be at most 10, the difference in strikes.

At expiry, the position pays 0 above 110, $110 - S$ between 100 and 110, and 10 below 100. So it never pays less than zero, and you were paid \$1 to take it.

\iqwrong{Thinking this can be explained by the smile. No smile can make a put at a higher strike cheaper: that breaks the rule in every model.}
\end{solution}

\begin{solution}{iq:iv2:options-parity-and-greeks:8}
A long call has positive delta, and a long put negative. Gamma and vega are positive for both, and theta is negative for both. The ones that are the same must be so, because by parity the two differ only by the stock and a bond, with no gamma, vega or (at zero rates) theta.

Parity says $C - P = S - K$ with zero rates. Taking derivatives, $\Delta_C - \Delta_P = 1$. All other derivatives are the same for both, since $S - K$ has none. So whatever gamma, vega or theta a call has, the put at that strike has as well.

\iqwrong{Saying a put has negative gamma because its delta is negative. Delta rises with the stock for both, from $-1$ toward 0 for a put.}
\end{solution}

\begin{solution}{iq:iv2:options-parity-and-greeks:9}
Your delta is 550 shares long. Sell 550 shares to make it zero.

Calls: $10 \times 100 \times 0.6 = 600$. Short puts: $-5 \times 100 \times (-0.3) = +150$. Stock: $-200$. And $600 + 150 - 200 = 550$.

\iqwrong{Forgetting the double negative on the short put. A short put gains when the stock rises, so it is long delta.}
\end{solution}

\begin{solution}{iq:iv2:options-parity-and-greeks:10}
You sell 250 shares to hedge. After a \$1 rise, the delta is 0.52, so the calls are long 260 shares, and you sell 10 more.

$5 \times 100 \times 0.5 = 250$. The new delta is $0.5 + 0.02$. So $5 \times 100 \times 0.52 = 260$. The hedge sells more as the stock rises, and buys more as it falls. Being long gamma means you sell high and buy low.

\iqwrong{Hedging once and thinking it is finished. Delta moves with the price, and gamma measures how fast.}
\end{solution}

\begin{solution}{iq:iv2:options-parity-and-greeks:11}
With a \$2 move, you make $\tfrac12 \times 1000 \times 4 - 1500 = +\$500$. With a \$1 move, you lose \$1,000. You break even at a move of about \$1.73.

Gamma's gain is $\tfrac12\Gamma m^2$. For the break-even, $500m^2 = 1500$, so $m = \sqrt 3$.

\iqwrong{Forgetting the half. It is there because the delta you earn on rises from zero to its full size over the move.}
\end{solution}

\begin{solution}{iq:iv2:options-parity-and-greeks:12}
You make \$1,100. A single move of 3 points on the whole vega would say \$1,500.

$300 \times 3 + 200 \times 1 = 1{,}100$. The surfaces do not move in parallel, as short-dated volatility moves more, so you need vega by expiry bucket.

\iqwrong{Quoting one total vega number for a book. A book can have zero vega and still lose heavily if the front and far months move differently.}
\end{solution}

\begin{solution}{iq:iv2:options-parity-and-greeks:13}
The call gains about \$0.47, and the put loses about \$0.53. A higher rate lowers the present value of the strike. The call owner pays that strike, and so gains, while the put owner receives it, and so loses.

Rho\index{rho} for a call is $KTe^{-rT}\Phi(d_2)$ per 1.00 of rate, here $100 \times \Phi(-0.1) \approx 46$, or 0.46 per point at a zero rate. For a put, it is $-100 \times \Phi(0.1) \approx -54$. Over a full point, the call's rho grows a little, and that gives 0.47 and $-0.53$. Parity gives the same total: a point on the present value of 100 is about 1.

\iqwrong{Thinking rates move stocks and options the same way. The stock is held fixed here. In practice you cannot hold everything fixed.}
\end{solution}

\begin{solution}{iq:iv2:options-parity-and-greeks:14}
An American call is worth at least the European's bound $S - Ke^{-rT}$, and that is above $S - K$, what exercise gives you. So selling it is always better than exercising it. For a put, exercising gets you the strike now, and the same bound runs the wrong way.

For a European call, the lower bound is $S - Ke^{-rT}$. With positive $r$ it exceeds $S - K$ by the interest on the strike. An American call is worth at least the European one. So exercise gives less than the option's price.

For a put, the bound is $Ke^{-rT} - S$, below the $K - S$ of exercise. So if the put is deep enough in the money, exercising earlier gets you the strike and its interest. Also, a call's insurance value disappears for nothing if you exercise it.

\iqwrong{``You should exercise if it is deep in the money.'' You should sell it, or hedge it, and not exercise it, as long as there is no dividend.}
\end{solution}

\begin{solution}{iq:iv2:options-parity-and-greeks:15}
The difference should be $100 - 95.12 = 4.88$, and it is 3. So the call is cheap against the put. Buy the call, sell the put, short the stock, and lend the rest. That gives you about \$1.88 today.

Buying the call and selling the put costs 3. Shorting the stock brings in 100. You lend 95.12 of the 97, and you keep 1.88. At expiry, the bond pays 100. The call and put together let you buy the stock for 100, however it ends, and you use it to close your short. So everything nets to zero, whatever the stock price.

This is called a reversal\index{reversal}, a conversion run the other way round, and it leaves you with no risk on the stock.

\iqwrong{Trading only the options. Without the stock and the bond, you own a view, and not an arbitrage.}
\end{solution}

\begin{solution}{iq:iv2:options-parity-and-greeks:16}
The fee takes the gap away. A gap of 1.88 on a \$100 stock is about 1.88\% for a year, so a 2\% fee more than covers it. The gap therefore implies a borrow of about 1.9\%. So it tells you what the short costs, and it is not an easy profit.

Doing the trade requires you to short the stock. Paying 2\% on 100 for a year costs about 2, which is more than the 1.88 you locked in.

A stock that is hard to borrow has a forward below the usual carry, and options show it through parity. Option markets are often the most accurate quote for a stock's borrow.

\iqwrong{Telling an interviewer you have found a mispricing. Ask first: is there a dividend, what is the borrow, and can the short be recalled?}
\end{solution}

\begin{solution}{iq:iv2:options-parity-and-greeks:17}
The limits are $S - K = 0$ and $S - Ke^{-rT} \approx 4.88$. So $0 \leq C - P \leq 4.88$.

For the upper limit, $C \leq S - Ke^{-rT} + P$. Since an American call on a no-dividend stock is worth the same as the European, the European parity gives $C = S - Ke^{-rT} + P_E$, and $P_E \leq P$. So $C - P \leq S - Ke^{-rT}$.

For the lower limit, buy the call, hold cash $K$, sell the put and short the stock. If the put is exercised against you at any time, you pay $K$ from the cash, receive a share and close the short, and you still have the call and the interest. If it is never exercised, then at expiry the call and the short put together pay $S_T - K$, and that with the short share and the cash leaves the interest. So the end value is never negative, and the position cannot cost less than zero: $C + K - P - S \geq 0$.

\iqwrong{Using equality. A short American put may be exercised against you at any time, so you cannot make the parity trade exact, and you get a range.}
\end{solution}

\begin{solution}{iq:iv2:options-parity-and-greeks:18}
It pays exactly \$20, wherever the stock finishes. So it is a one-year zero-coupon bond. At \$18.70 it implies about 6.72\%, so you borrow at 5\% and buy it.

The first two legs make a long forward at 90, and the last two a short forward at 110. Together they pay $(S - 90) - (S - 110) = 20$. So $18.70e^{r} = 20$ gives $r = \ln(20/18.70) \approx 6.72\%$. At 5\%, the box's fair value is $20e^{-0.05} \approx 19.02$.

This is a box spread\index{box spread}: a way of lending or borrowing through options. So it is a way to measure the rate implied in options.

\iqwrong{Treating it as free. It relies on your being able to hold the position to expiry, and American options, margin and the credit of the clearing house all come into it.}
\end{solution}

\begin{solution}{iq:iv2:options-parity-and-greeks:19}
Yes, on a stock with no dividends it is an arbitrage. The longer call must be worth at least the shorter one. With a large dividend in between, it is not: holders of the longer call miss out on the dividend. To trade, you sell the short one at \$4, buy the long one at \$3.50, and pocket \$0.50.

At the first expiry, if the short call is in the money, you owe $S - 100$. Your long call has two months left, and it is worth at least $S - 100e^{-r\tau}$, which is more than $S - 100$. So you can sell it, and you always cover it.

With a dividend between the expiries, the price drop lowers the longer call, and its value can be lower.

\iqwrong{Saying a longer option always costs more. It needs the assumption of no dividend in between, and, for a European put, it can also fail with high rates.}
\end{solution}

\begin{solution}{iq:iv2:options-parity-and-greeks:20}
Gamma is about 0.0198 with a year left and 0.317 with one day left, about 16 times as much. With one day left, gamma halves at about \$1.48 from the strike. So short-dated short options carry a big and rapidly changing risk.

$\Gamma = \varphi(d_1)/(S\sigma\sqrt T)$. With a year: $\varphi(0.1)/20$. With a day: $\sigma\sqrt T = 0.2/\sqrt{252} \approx 0.0126$, so $\Gamma \approx 0.399/1.26 \approx 0.317$. Gamma is $\varphi$ of the standardised distance, so it halves when $d = \sqrt{2\ln 2} \approx 1.18$, that is, at about 1.18 daily standard deviations, or \$1.48.

So a short gamma position near a strike at the end of its life is a big exposure. Across a \$3 range centred on the strike, the delta goes from about 0.12 to about 0.88.

\iqwrong{Thinking risk falls as an option nears expiry because the option becomes cheap. Its value falls, and its gamma, near the strike, goes up.}
\end{solution}

\begin{solution}{iq:iv2:options-parity-and-greeks:21}
It is worth about \$45.13. It gains about 1.3 cents with one more day passing, or about \$4.74 a year. Theta can be positive in this way when a European put is deep in the money with a positive rate, or when a deep in-the-money call has a high dividend yield.

The put is almost certain to pay 100 less the stock, and so it is nearly the bond: $100e^{-0.05} - 50 \approx 45.12$. As time passes, the discounted strike grows, by $rKe^{-rT} \approx 4.76$ a year, and that is about the theta.

\iqwrong{``Theta is always negative.'' That is true for an option at zero rates. A deep in-the-money European put is mostly a bond, and it accrues like one.}
\end{solution}

\begin{solution}{iq:iv2:options-parity-and-greeks:22}
The vegas are about 0.19, 0.40 and 0.29 per point, at strikes of 80, 100 and 120. They are not symmetric around 100, for two reasons: vega peaks at $d_1 = 0$, above the stock, and a 20-point distance is larger in log terms below than above.

Vega is $S\varphi(d_1)\sqrt T$. The $d_1$ values are about 1.22, 0.10 and $-0.81$. The peak, $d_1 = 0$, is at the strike $Se^{\sigma^2T/2} \approx 102$. And the 120 strike is closer in logs: $\ln(120/100) = 0.18$, against $\ln(80/100) = -0.22$.

So a strike 20\% above has more vega than one 20\% below. This matters in setting up a vega-neutral trade.

\iqwrong{Assuming that equal dollar distances have equal Greeks. The normal distribution lives in log terms.}
\end{solution}

\begin{solution}{iq:iv2:options-parity-and-greeks:23}
The delta is $\Phi(0.2) \approx 0.579$, and the chance is $\Phi(-0.2) \approx 0.421$. The trader is off by about 0.16. The claim is right only when total volatility is small.

Here, $d_1 = \sigma\sqrt T/2 = 0.2$ and $d_2 = -0.2$. The difference between them is $\sigma\sqrt T$. So the probability that the option finishes in the money is $\Phi(d_2)$, under the risk-neutral measure, and the delta is $\Phi(d_1)$.

Delta is the chance under the measure that has the stock as numeraire. That measure gives high prices more weight, which is why it is higher.

\iqwrong{Saying that delta is the probability of finishing in the money. They are close for short-dated, low-vol options, and that is where the habit comes from.}
\end{solution}

\begin{solution}{iq:iv2:options-parity-and-greeks:24}
The put's delta is about $-0.43$. With no dividend yield, it would be $-0.45$.

With dividend yield $q$, parity is $C - P = Se^{-qT} - Ke^{-rT}$, and differentiating gives $\Delta_C - \Delta_P = e^{-qT}$. So $\Delta_P = 0.55 - e^{-0.02} \approx 0.55 - 0.9802 = -0.4302$.

\iqwrong{Using $\Delta_P = \Delta_C - 1$ in all cases. That is the no-dividend case, and it is off by 2\% here.}
\end{solution}

\begin{solution}{iq:iv2:options-parity-and-greeks:25}
You sell 0.12 shares at 103, and you buy them back at 100, for a gain of \$0.36 per straddle. The formula says $2 \times \tfrac12 \times 0.04 \times 9 = 0.36$. A day's theta of $-0.36$ would break even.

At 100, delta is zero. At 103, delta is $0.04 \times 3 = 0.12$, so you sell 0.12 at 103. Back at 100, the delta is zero again, so you buy 0.12 at 100. That makes $0.12 \times 3 = 0.36$. So the gain comes from selling high and buying low. This is gamma scalping\index{gamma scalping}.

\iqwrong{Thinking that an unchanged stock means no profit on the day. The hedger's gain depends on how far the stock moved in between.}
\end{solution}

\begin{solution}{iq:iv2:options-parity-and-greeks:26}
You should have made \$50: $+400$ from delta, $+100$ from gamma, $-300$ from vega and $-150$ from theta. You made \$120, so \$70 is unexplained.

Delta: $200 \times 2$. Gamma: $\tfrac12 \times 50 \times 4$. Vega: $300 \times (-1)$. Theta: $-150$.

The \$70 could be a change in skew, a move in vega with spot (vanna), an intraday path if you rehedged, dividends, a data error, or trades you did today at a good price. Before you trust the model, find out which. A persistent unexplained amount is a warning sign.

\iqwrong{Ignoring a small unexplained amount because the total is close. Errors may cancel, and an unexplained profit is as much of a warning as a loss.}
\end{solution}

\begin{solution}{iq:iv2:options-parity-and-greeks:27}
At expiry, you lose \$9.50 at 80, you make \$0.50 at 100, and you make \$10.50 at 120. The delta is about 0.62 and the vega about 0.06 per point. You gain about \$0.68 from a one-point rise in the 110 implied volatility and a one-point fall in the 90.

The net credit is 0.50. At 80, the put costs you 10. At 120, the call pays 10. The delta is $0.35 + 0.27$, and the vega $0.37 - 0.31$. A rise in the call's volatility makes 0.37, and a fall in the short put's makes 0.31. So you own a risk reversal\index{risk reversal}, which is mostly a bet on the slope of the smile, and on direction.

\iqwrong{Calling this free because it costs nothing or pays. A credit does not mean no risk, as you have all of the downside below 90.}
\end{solution}

\begin{solution}{iq:iv2:options-parity-and-greeks:28}
Your largest gain and loss are both \$10. The position is the same as buying a 90-110 call spread and holding \$90 in cash. That call spread must cost \$10.

At expiry, your holding is worth $\min(\max(S, 90), 110)$. That equals $90 + \max(S - 90, 0) - \max(S - 110, 0)$. By parity with zero rates, $S + P_{90} = C_{90} + 90$. So $S + P_{90} - C_{110} = 90 + C_{90} - C_{110}$. The cost on the left is $100 + 3 - 3 = 100$, so the spread costs 10.

\iqwrong{Calling a collar\index{collar} ``free insurance''. You paid for the protection by selling your gain above 110.}
\end{solution}

\begin{solution}{iq:iv2:options-parity-and-greeks:29}
Your profit is $-\$12$ at 85, $+\$8$ at 105 and $+\$8$ at 120. It is the same as selling a 105 put, and that must trade at \$8.

$S + (-\text{call}) = \min(S, 105)$ at expiry, plus the premium, less the 100 you paid. By parity with zero rates, $P = C - S + K = 3 - 100 + 105 = 8$. A short put at 8 pays $8 - 20 = -12$ at 85, and 8 above 105.

So covered-call selling is selling a put. You keep nearly all the downside, and you give up the upside, for a modest premium.

\iqwrong{``Income.'' The premium is pay for taking risk. The position is a short put, and in a big fall it is much like owning stock.}
\end{solution}

\begin{solution}{iq:iv2:options-parity-and-greeks:30}
The call less the put is about \$4.88: $(100 - 95)e^{-0.025}$.

Long a call and short a put at the same strike is a forward agreement to buy at 95. The future's value at expiry is $F$, so the payoff is $F_T - 95$. A futures contract costs nothing to enter, so its fair value today is the present value of $F - K$. So $C - P = (F - K)e^{-rT}$.

The stock version has $S - Ke^{-rT}$: the stock has to be paid for today, while a future does not. If you write $F = Se^{rT}$, the two are the same.

\iqwrong{Using $S - Ke^{-rT}$ with the future's price in place of $S$. That gives $100 - 95e^{-0.025} \approx 7.35$, and it counts a financing cost you are not paying.}
\end{solution}

\begin{solution}{iq:iv2:options-parity-and-greeks:31}
It is worth about \$0.46, which is $\Phi(d_2)$ with $d_2 = -0.1$. It is the limit of a tight call spread: $-\partial C/\partial K$. Near expiry and near the strike, its delta becomes huge, and the hedge jumps.

A spread paying 1 above $K + \varepsilon$, made of $1/(2\varepsilon)$ calls at $K - \varepsilon$ and $K + \varepsilon$, tends to the digital. Its price is $-\partial_K C$. In Black-Scholes, this is $\Phi(d_2)$. A numerical check at $\varepsilon = 0.5$ gives the same to three decimals.

Delta is $\varphi(d_2)/(S\sigma\sqrt T)$, and it grows without limit as $T \to 0$, with the stock at the strike.

\iqwrong{Using a call's delta, $\Phi(d_1)$, as the price. That would be 0.54 here. A digital in cash is valued by the risk-neutral probability.}
\end{solution}

\begin{solution}{iq:iv2:options-parity-and-greeks:32}
Its gamma is positive near and below 100, negative near and above 110, and it crosses zero between them, near the middle of the range. This holds because it is long gamma at the lower strike and short at the higher, and each option's gamma peaks near its own strike.

In Black-Scholes, the two gammas are equal where the two $d_1$ are equal and opposite, at $S = \sqrt{K_1K_2}\,e^{-\sigma^2T/2}$. With a month left at 20\% volatility, that is about 104.7.

So as the stock rises through the range, your delta first rises, and then falls. Rehedging, you sell into a rise near 100, and you buy into a rise near 110. The spread pays theta near the lower strike, and it collects theta near the upper.

\iqwrong{Treating a spread's Greeks as fixed. Their sign depends on where the stock is, so the hedge can reverse during a day.}
\end{solution}

\begin{solution}{iq:iv2:options-parity-and-greeks:33}
The call at 25\% is about \$9.95, and the put at 27\% is about \$10.74. So they are about \$0.79 from parity, which is one way of saying that the two implied volatilities are 2 points apart. The implied borrow is about 0.8\% a year.

With zero rates, parity has $C = P$ at the strike. A vega of 0.40 per point on a 2-point gap gives about \$0.80. Lower the forward by \$0.79, and both options fit one volatility. So what you see as a volatility gap is really a wrong forward.

So most systems solve for one implied forward from the call and put near the money. They then use one volatility per strike, usually the out-of-the-money option's.

\iqwrong{Trying to arbitrage the two volatilities without asking about the forward. The arbitrage needs the stock: if it is costly to short, the gap can stay.}
\end{solution}

\begin{solution}{iq:iv2:options-parity-and-greeks:34}
Yes, exercise now. Exercising gets you 99 today, while a European put can never be worth more than the strike's present value, $100e^{-0.05} \approx 95.12$. At a zero rate, you would not exercise it, because there is no interest to earn on the strike, and you would give up the option's remaining value.

Holding the put to expiry pays at most 100 in a year, and that is worth 95.12 today. Taking 99 now and investing it gives about 104.08 in a year. So early exercise beats the best case for holding.

\iqwrong{Using the call result for a put. For a call, delay is a benefit: you keep the strike. For a put, delay costs you: you are waiting to get paid.}
\end{solution}

\begin{solution}{iq:iv2:options-parity-and-greeks:35}
With a year left, you hold about 19,850 shares. With one day left, you hold about 316,600 shares, on a \$1 million payoff. A dealer prices it as a call spread that always pays at least the digital. That caps your delta and lets you hedge with ordinary options.

The hedge is $\$1{,}000{,}000 \times \varphi(d_2)/(S\sigma\sqrt T)$. With a year: $1{,}000{,}000 \times 0.397/20 \approx 19{,}850$. With a day: about 316,600. So on the day, a \$1 stock move gives a big change, and you must trade a huge number of shares right around the strike.

To reduce the risk, you hold a call spread of width $w$ against it: long $1{,}000{,}000/w$ calls at $K - w$ and short as many at $K$. It pays the full amount above $K$, so it covers what you owe, and a little more just below. With $w = 2$, its delta never exceeds 500,000 shares, however close the expiry, where the digital's delta has no limit. The extra premium is what you charge for this comfort.

\iqwrong{Hedging a digital as you hedge a vanilla and trusting the model's delta. The model's delta is enormous near the strike, and in a real market you cannot trade it.}
\end{solution}

\begin{solution}{iq:iv2:options-parity-and-greeks:36}
Measured from the close, if you are assigned on all of it, you lose nothing. If you are not assigned at all, you lose \$200,000. This is pin risk\index{pin risk}. You should have closed the position, or at least reduced it, before the expiry.

At the close, each share less its short call was worth $100.02 - 0.02 = 100$.

Assigned: you deliver your 100,000 shares at \$100, which is that same value, and you have no further stock risk.

Not assigned: the calls expire, and you still hold 100,000 shares, now at \$98. That is $-\$2$ a share, or \$200,000.

And you should expect the bad case. A holder who sees \$98 before the deadline will not exercise, and one who sees \$102 will, so the choice runs against you. Close the hedge and the options at the end of the day, even at a small cost, so that you do not carry an unknown position overnight.

\iqwrong{Calling the position delta-neutral. On the final day, your delta is either 0 or 100,000 shares, and you find out which after the market has moved.}
\end{solution}

\begin{solution}{iq:iv2:options-parity-and-greeks:37}
Yes, exercise. You gain about \$0.36 over waiting. The rule is to exercise when the dividend is more than $K(1 - e^{-r\tau})$ plus the put's value. Here $2 > 1.24 + 0.40$.

Exercising now gives $110 - 100 = 10$. If you wait, the stock falls to 108. By parity with no more dividends, the call is worth $P + (S - D) - Ke^{-r\tau} = 0.40 + 108 - 98.76 = 9.64$. So the gap is 0.36.

In general, you compare $S - K$ with $P + S - D - Ke^{-r\tau}$. So exercise if $D > K(1 - e^{-r\tau}) + P$. You receive the dividend, and you give up two things: the interest on the strike and the insurance a live option provides.

So exercise is decided just before the ex-date. A market maker has to watch every American option before each dividend.

\iqwrong{Exercising because the call is deep in the money. Depth alone never justifies it, and only the dividend does.}
\end{solution}

\begin{solution}{iq:iv2:options-parity-and-greeks:38}
The European put is worth about \$6.29, and the American about \$7.99. You exercise early at the down node after one period, when the stock is at 80, and the premium is about \$1.70.

The risk-neutral up-probability is $(1.05 - 0.8)/(1.2 - 0.8) = 0.625$. The final payoffs are 0 at 144, 4 at 96, and 36 at 64.

At the up node (120), holding is worth $(0.625 \times 0 + 0.375 \times 4)/1.05 \approx 1.43$, and exercising gives 0. At the down node (80), holding is worth $(0.625 \times 4 + 0.375 \times 36)/1.05 \approx 15.24$, and exercising gives 20. So you exercise there.

The European value is $(0.625 \times 1.43 + 0.375 \times 15.24)/1.05 \approx 6.29$. The American value is $(0.625 \times 1.43 + 0.375 \times 20)/1.05 \approx 7.99$.

Pricing American options on a tree\index{binomial tree} comes down to this comparison at each node.

\iqwrong{Comparing exercise with the option's final value. You have to compare with the discounted expected value at that node.}
\end{solution}

\begin{solution}{iq:iv2:options-parity-and-greeks:39}
Near 100, a day's gain is about \$0.092. Near 120, it is about \$0.011. So the same realised volatility gives very different profits, and the first is more than 8 times the second. Realised volatility alone does not measure your profit.

A move of 2\% at 100 is \$2, and $\tfrac12 \times 0.046 \times 4 = 0.092$. At 120 it is \$2.40, and $\tfrac12 \times 0.0038 \times 5.76 \approx 0.011$.

Your profit along a path is $\sum\tfrac12\Gamma_tS_t^2(r_t^2 - \sigma_i^2\,dt)$, gamma-weighted. So days when the stock is near the strike count most. So realised volatility on those days is what matters. A variance swap\index{variance swap} is designed to remove this dependence, by holding a constant dollar gamma.

\iqwrong{Saying that you win whenever realised volatility exceeds implied. You can see a realised above implied and still lose if the big moves came when you had little gamma.}
\end{solution}

\begin{solution}{iq:iv2:options-parity-and-greeks:40}
You buy about 1.135 strangles, you collect about \$6.99, and you make about \$0.10 if volatility rises to 25\%, and about \$0.16 if it falls to 15\%. You are long volga\index{volga}: your vega grows in the direction of the move.

The vega of the straddle is about 79.4 per 1.00, and the strangle's is about 69.9, so the ratio is 1.135. At 20\%, the straddle costs about 15.93, and 1.135 strangles cost about 8.95. So you collect about 6.99.

An at-the-money option's vega hardly moves with volatility, as its price is nearly linear in volatility. An out-of-the-money option's vega grows with volatility. So this position has zero vega at first, and a positive second derivative. That is, it gains from a large move in implied volatility, either way.

\iqwrong{Thinking zero vega means no volatility exposure. It means no first-order exposure at this level, and the second-order one is exactly what the position holds.}
\iqwhy{Volga is what a dealer who is short the wings is paid to carry, and that is one reason wings trade at a higher implied volatility.}
\end{solution}

\begin{solution}{iq:iv2:options-parity-and-greeks:41}
The standard deviation is about \$0.44, about 5.6\% of the \$7.97 premium. Hedging four times as often halves it.

The approximation is $\sqrt{\pi/4}\,\sigma\,\mathcal V/\sqrt N$, where $\mathcal V$ is the vega per 1.00 of volatility: $0.886 \times 0.2 \times 39.7/\sqrt{252} \approx 0.443$. A simulation of 20,000 paths gives about 0.44.

The error is proportional to $1/\sqrt N$. Each step leaves a residual proportional to gamma times the square of the move less its expected value, and these errors add as independent variances.

So a dealer has to charge for this risk. It is real even if the true volatility is exactly the implied. Hedging more often lowers it, but at a higher cost.

\iqwrong{Saying that the option is ``risk-free once hedged''. That holds only for continuous hedging, and there is always a residual, which falls slowly with frequency.}
\end{solution}

\begin{solution}{iq:iv2:options-parity-and-greeks:42}
Raise the volatility from 20\% to about 21.2\%, which adds about \$0.49 to the premium. A simulation of the rebalancing costs gives about \$0.50.

Leland's\index{Leland's rule} rule is $\hat\sigma^2 = \sigma^2(1 + A)$, where $A = \sqrt{2/\pi}\,k/(\sigma\sqrt{\Delta t})$ and $k$ is the round-trip cost. Here $\sigma\sqrt{\Delta t} = 0.0126$, and $A = 0.798 \times 0.002/0.0126 \approx 0.127$. So $\hat\sigma = 0.2\sqrt{1.127} \approx 21.2\%$. With a vega of 39.7 per 1.00, that is about $0.0123 \times 39.7 \approx 0.49$.

The expected number of shares you trade per step is $\Gamma S\sigma\sqrt{\Delta t}\sqrt{2/\pi}$, since the expected absolute move is that fraction of a standard deviation. The cost of those trades, charged to gamma, looks like extra volatility. The simulated figure leaves out the first purchase of about 54 shares per 100, which adds about \$0.05.

A seller raises the volatility, and a buyer lowers it, because the costs go against you on either side.

\iqwrong{Assuming more hedging always helps. As $\Delta t$ falls, $A$ grows, and at some point your costs outweigh the risk reduction. So choosing the frequency is itself the problem.}
\end{solution}

\begin{solution}{iq:iv2:options-parity-and-greeks:43}
You need about 3.45 one-month lots, and the net gamma is about $+0.218$ per dollar. If volatility moves in proportion, you only need about 1 lot.

The vegas are about 39.7 and 11.5 per 1.00. So $39.7/11.5 \approx 3.45$. The gammas are about 0.0198 and 0.0691. The net is $3.45 \times 0.0691 - 0.0198 \approx 0.218$.

In practice, short-dated implied volatility moves more. If a one-month move is $\sqrt{12}$ times that of the one-year, one lot of one-month vega has $11.5 \times 3.46 \approx 39.8$ of one-year-equivalent vega. So one lot of each cancels. This is the idea behind a weighted vega\index{weighted vega}.

\iqwrong{Treating a vega-flat calendar as flat. It is flat to a parallel move only, and it is a very long gamma position.}
\end{solution}

\begin{solution}{iq:iv2:options-parity-and-greeks:44}
Not always. The value change per point is about $+0.0070$ at 90, $-0.0020$ at 100 and $-0.0107$ at 110. The sign changes at about 98.02, not at 100.

The value is $\Phi(d_2)$, so its derivative in $\sigma$ is $\varphi(d_2)\,\partial_\sigma d_2 = -\varphi(d_2)d_1/\sigma$. That is zero when $d_1 = 0$, so $S = Ke^{-\sigma^2T/2} = 100e^{-0.02} \approx 98.02$.

Below that, more volatility gives a better chance of reaching the strike. Above it, more volatility gives a greater chance of falling back. An ordinary call has a vega that is positive everywhere, since it is never hurt by more volatility.

\iqwrong{Assuming that all options gain from volatility. Payoffs that are not convex in the stock can lose from it.}
\end{solution}

\begin{solution}{iq:iv2:options-parity-and-greeks:45}
If they are long gamma and the stock rises \$1, they sell 1 million shares, 20\% of a day's volume, and if it falls, they buy. If they are short, they buy into the rise and sell into a fall, and that adds to the move. So long gamma tends to dampen moves and pin the price, and short gamma tends to amplify them.

The estimate is unreliable. It assumes the dealers are on one side of all customer trades, and it does not know how they have hedged. Nor do open-interest numbers tell you who bought and who sold. So check the claim against the data on actual realised volatility, and treat it as a weak prior.

\iqwrong{Stating it as a fact of the market. It is a mechanism with evidence, but it rests on an estimate you cannot see directly.}
\end{solution}

\begin{solution}{iq:iv2:options-parity-and-greeks:46}
Differentiate $C(\lambda S, \lambda K) = \lambda C(S, K)$ in $\lambda$ at 1. This gives $C = S\partial_S C + K\partial_K C$. So the delta is $(C - K\partial_K C)/S$. From the quotes, $\partial_K C \approx (7.52 - 8.44)/2 = -0.46$. So the delta is $(7.97 + 100 \times 0.46)/100 = 0.5397$, against the Black-Scholes $\Phi(0.1) \approx 0.5398$.

This assumes that the call's price depends only on the ratio of stock to strike, so the smile moves with the stock: sticky moneyness\index{sticky moneyness}. Then the delta can be read from a call's slope in strike, and needs no model. If instead the smile stays fixed by strike (sticky strike), the delta is a different number, with skew adding to or subtracting from it.

\iqwrong{Thinking there is one correct delta. The model that produces the delta includes a belief about how the smile moves.}
\end{solution}

\begin{solution}{iq:iv2:options-parity-and-greeks:47}
All prices fall with the strike, and no spread exceeds its width. The convexity condition breaks at 100, 110 and 120. So buy the 100 and the 120 and sell two of the 110. You are paid \$1, and your worst outcome is that \$1 profit, and your best is \$11. The table implies chances of about 60\%, 25\% and 35\%, and the rise from 25\% to 35\% is impossible.

The butterfly costs $8 - 2 \times 5.5 + 2 = -1$. Its payoff is 0 at or below 100 and above 120, and it reaches 10 at 110. So $-(-1) + 0 = 1$ at worst, and $1 + 10 = 11$ at best.

Slopes: $(14 - 8)/10 = 0.60$, $(8 - 5.5)/10 = 0.25$, $(5.5 - 2)/10 = 0.35$. A call spread's value per unit of width is about the chance of finishing above its middle, and that must fall with the strike. So this is the same breach as the butterfly's.

\iqwrong{Testing only that prices fall with the strike. Convexity is the condition that fails most often.}
\end{solution}

\begin{solution}{iq:iv2:options-parity-and-greeks:48}
Any twice-differentiable $f$ satisfies
\[
f(S) = f(K_0) + f'(K_0)(S - K_0) + \int_0^{K_0}f^{\prime\prime}(K)(K - S)^+\,dK + \int_{K_0}^\infty f^{\prime\prime}(K)(S - K)^+\,dK.
\]
For $S_T^2$, $f^{\prime\prime} = 2$ for every strike. So you hold 2 of every put below 100 and 2 of every call above it. With zero rates, the strip comes to about 408.11. With the bond of 10,000 and the forward's zero value, the total is 10,408.11, which is exactly $S^2e^{\sigma^2T}$.

The formula comes from Taylor's theorem with an integral remainder. Integrate by parts, and the second-derivative terms are options. This is the Carr-Madan\index{Carr-Madan replication} static replication.

For $f = S^2$: $f(K_0) = 10{,}000$, and $f'(K_0)(S - K_0)$ is a forward with a value of zero. The options are weighted $f^{\prime\prime}(K)\,dK = 2\,dK$.

For the check, Black-Scholes gives $E[S_T^2] = S_0^2e^{\sigma^2T} = 10{,}000e^{0.04} \approx 10{,}408.11$. Numerical integration of the strip agrees.

So any European payoff has a model-free value, given the option prices, with no model of the path. It is the basis of the variance swap's price, using $f = \ln$.

\iqwrong{Thinking you need a model of volatility to price a smooth payoff. The smile prices it, and it is only a matter of how many strikes you can trade.}
\end{solution}

\begin{solution}{iq:iv2:options-parity-and-greeks:49}
Your expected profit is about \$1.99, the difference between the option at 30\% and at 20\%. Hedging with implied deltas gives a profit along a path of $\int\tfrac12\Gamma_iS^2(\sigma_i^2 - \sigma_r^2)\,dt$, with a daily-hedge mean of about 1.99 and a standard deviation of about 0.74. Hedging with realised deltas gives a sure profit of $V(\sigma_i) - V(\sigma_r)$ under continuous hedging, and in daily hedging, about 1.99 with a standard deviation of about 0.44. If you knew the 20\%, you would hedge with it. Most desks use the implied, since no one knows the realised, and the P\&L then accrues smoothly at a known rate.

Hedge at $\sigma_r$. Let $V_r(S, t)$ be the option's value at volatility $\sigma_r$. By Ito, its change is $\Delta_r\,dS$, with theta and gamma cancelling at that realised volatility. So the hedge offsets the option from the start, at a value of $V_r$, and you keep $V_i - V_r$ in every path. The 0.44 is only the discreteness.

Hedge at $\sigma_i$. Apply Ito to $V_i$ under the realised dynamics: $dV_i - \Delta_i\,dS = (\theta_i + \tfrac12\sigma_r^2S^2\Gamma_i)dt = \tfrac12\Gamma_iS^2(\sigma_r^2 - \sigma_i^2)dt$. For a short option, you keep the negative of that. Taking the expected value gives the same $V_i - V_r$.

\iqwrong{Saying the implied hedge is ``right'' because that is the market's volatility. It does not reduce variance, since the realised path is what drives your P\&L. Its advantage is in the timing of the P\&L.}
\end{solution}

\begin{solution}{iq:iv2:options-parity-and-greeks:50}
The exercise level is about \$71.43, and the put is worth about \$12.32 with the stock at 100. As the rate goes to 0, you never exercise, and the value tends to the strike.

A perpetual option has no time dependence, so the Black-Scholes equation becomes $\tfrac12\sigma^2S^2V^{\prime\prime} + rSV' - rV = 0$. Try $V = S^p$: $\tfrac12\sigma^2p(p - 1) + rp - r = 0$, with roots $p = 1$ and $p = -\gamma$, where $\gamma = 2r/\sigma^2 = 2.5$. The value must vanish for large $S$, so $V = AS^{-\gamma}$.

At the exercise level $S^*$, $V(S^*) = K - S^*$. Then for a level $L$, the value at $S$ is $(K - L)(S/L)^{-\gamma}$. To maximise over $L$: the derivative of $(K - L)L^\gamma$ is zero at $L = \gamma K/(1 + \gamma) = 71.43$. This choice is the same as smooth pasting: $V'(S^*) = -1$.

With the stock at 100: $28.57 \times (1.4)^{-2.5} \approx 12.32$.

As $r \to 0$, $\gamma \to 0$, so $L \to 0$, and $V \to K$. With no interest to earn on the cash, you wait. At zero drift, the stock tends to zero with probability one, so the put's value tends to the whole strike.

\iqwrong{Fixing the exercise level by setting the value equal to a European's value. The level comes from maximising over the choice of exercise, and that is the essence of American pricing.}
\iqfollow{What does a finite expiry do?}{Then no closed form exists. A tree or a finite-difference grid solves it, and the boundary rises with the passage of time toward the strike.}
\end{solution}
